\documentclass[aspectratio=169]{beamer}

% Use the UFS theme
\usetheme{UFS}

% Additional packages
\usepackage[utf8]{inputenc}
\usepackage[portuguese]{babel}
\usepackage{amsmath,amssymb}
\usepackage{graphicx}
\usepackage{booktabs}
\usepackage{listings}
\usepackage{clrscode3e}
\usepackage{xcolor}
\usepackage{tikz}
\usepackage{colortbl}
\usepackage{pgfplots}

\usetikzlibrary{shapes.geometric, arrows, positioning, matrix, fit, backgrounds, decorations.pathreplacing, shadows, patterns, automata, intersections, calc,decorations.pathmorphing, arrows.meta}

% --- Paleta de cores ---
\definecolor{nodeprev}{HTML}{8E44AD}
\definecolor{highlight}{RGB}{200,40,40}
\definecolor{bgcolor}{HTML}{FFFFFF}
\definecolor{linecolor}{HTML}{DDDDDD}
\definecolor{numcolor}{HTML}{999999}
\definecolor{kwcolor}{HTML}{0000FF}
\definecolor{strcolor}{HTML}{A31515}
\definecolor{cmtcolor}{HTML}{008000}
\definecolor{fncolor}{HTML}{795E26}
\definecolor{typecolor}{HTML}{267F99}
\definecolor{nodedummy}{HTML}{9B59B6}
\definecolor{nodedead}{HTML}{888888}
\definecolor{nodecurr}{HTML}{F39C12}
\definecolor{nodecolor}{HTML}{4A90D9}
\definecolor{nodehigh}{HTML}{E74C3C}
\definecolor{nodeinsert}{HTML}{27AE60}
\definecolor{arrowold}{HTML}{AAAAAA}
\definecolor{arrownew}{HTML}{E74C3C}
\definecolor{destaque}{HTML}{E8F5E9}
\definecolor{alerta}{HTML}{FFF3E0}

\lstdefinestyle{modernstyle}{
    backgroundcolor=\color{bgcolor},
    basicstyle=\ttfamily\footnotesize\color{black},
    keywordstyle=\color{kwcolor}\bfseries,
    stringstyle=\color{strcolor},
    commentstyle=\color{cmtcolor}\itshape,
    numberstyle=\tiny\color{numcolor},
    numbers=left,
    numbersep=12pt,
    xleftmargin=20pt,
    framexleftmargin=20pt,
    breaklines=true,
    breakatwhitespace=false,
    tabsize=4,
    keepspaces=true,
    showspaces=false,
    showstringspaces=false,
    showtabs=false,
    frame=single,
    rulecolor=\color{linecolor},
    framesep=4pt,
    captionpos=b,
    upquote=true,
    columns=fullflexible,
}
\lstset{style=modernstyle}

% --- Metadados ---
\title{Estruturas Lineares: Arrays Dinâmicos e Listas Encadeadas}
\subtitle{COMP0497 -- Algoritmos e Estrutura de Dados 1}
\author{Prof.\ Dr.\ André Yoshiaki Kashiwabara}
\date{}

% =====================================================================
%  ESTILOS TIKZ: Listas Simplesmente Ligadas
% =====================================================================
\tikzset{
    llnode/.style={
        draw, rectangle, minimum width=1.4cm, minimum height=0.7cm,
        rounded corners=2pt, font=\ttfamily\small,
        fill=nodecolor!15, draw=nodecolor!70, thick
    },
    llnode high/.style={
        llnode, fill=nodehigh!15, draw=nodehigh!70
    },
    llnode insert/.style={
        llnode, fill=nodeinsert!15, draw=nodeinsert!70
    },
    llnode ghost/.style={
        llnode, fill=black!5, draw=black!25, text=black!35
    },
    ptr/.style={
        -Stealth, thick, color=nodecolor!80
    },
    ptr old/.style={
        -Stealth, thick, color=arrowold, dashed
    },
    ptr new/.style={
        -Stealth, very thick, color=arrownew
    },
    lbl/.style={
        font=\ttfamily\scriptsize\bfseries, text=nodehigh!80!black
    },
    lbl green/.style={
        font=\ttfamily\scriptsize\bfseries, text=nodeinsert!80!black
    },
    step label/.style={
        font=\sffamily\scriptsize\itshape, text=black!60
    },
    null node/.style={
        font=\ttfamily\small\bfseries, text=black!40
    },
}

% =====================================================================
%  ESTILOS TIKZ: Convenções de Lista
% =====================================================================
\tikzset{
    lnode/.style={
        draw, rectangle split, rectangle split parts=2,
        rectangle split horizontal, minimum height=0.55cm,
        thick, rounded corners=1pt,
        fill=nodecolor!12, draw=nodecolor!55,
        font=\ttfamily\scriptsize
    },
    lnode hi/.style={
        lnode, fill=nodehigh!12, draw=nodehigh!55
    },
    lnode new/.style={
        lnode, fill=nodeinsert!12, draw=nodeinsert!55
    },
    lnode sent/.style={
        lnode, fill=nodedummy!10, draw=nodedummy!50
    },
    larrow/.style={-stealth', thick, color=nodecolor!60},
    harrow/.style={-stealth', very thick, color=nodehigh!60},
    head ptr/.style={font=\ttfamily\tiny\bfseries, text=nodehigh!80!black},
    null text/.style={font=\ttfamily\scriptsize\bfseries, text=black!35},
    conv title/.style={font=\sffamily\scriptsize\bfseries, text=black!70},
    conv subtitle/.style={font=\sffamily\tiny\itshape, text=black!50},
}

% =====================================================================
%  ESTILOS TIKZ: Josephus
% =====================================================================
\tikzset{
    jnode/.style={
        draw, circle, minimum size=0.85cm, thick,
        fill=nodecolor!15, draw=nodecolor!70,
        font=\sffamily\bfseries\small
    },
    jnode curr/.style={
        jnode, fill=nodecurr!25, draw=nodecurr!80
    },
    jnode target/.style={
        jnode, fill=nodehigh!25, draw=nodehigh!80
    },
    jnode dead/.style={
        jnode, fill=black!8, draw=black!25, text=black!30
    },
    jnode survivor/.style={
        jnode, fill=nodeinsert!25, draw=nodeinsert!80,
        minimum size=1cm, font=\sffamily\bfseries
    },
    jarrow/.style={
        -stealth', thick, color=nodecolor!60
    },
    jarrow curr/.style={
        -stealth', very thick, color=nodecurr!80
    },
    jarrow skip/.style={
        -stealth', very thick, color=nodehigh!70, dashed
    },
    jlbl/.style={
        font=\sffamily\tiny, text=black!50
    },
}

% =====================================================================
%  ESTILOS TIKZ: Listas Duplamente Ligadas
% =====================================================================
\tikzset{
    dnode/.style={
        draw, rectangle split, rectangle split parts=3,
        rectangle split horizontal, minimum height=0.65cm,
        thick, rounded corners=1pt,
        fill=nodecolor!10, draw=nodecolor!60,
        font=\ttfamily\scriptsize,
        rectangle split part fill={nodeprev!8, nodecolor!8, nodecolor!8}
    },
    dnode hi/.style={
        dnode, fill=nodehigh!10, draw=nodehigh!60,
        rectangle split part fill={nodeprev!12, nodehigh!12, nodehigh!12}
    },
    dnode new/.style={
        dnode, fill=nodeinsert!10, draw=nodeinsert!60,
        rectangle split part fill={nodeinsert!10, nodeinsert!10, nodeinsert!10}
    },
    dnode ghost/.style={
        dnode, fill=black!4, draw=black!20, text=black!30,
        rectangle split part fill={black!4, black!4, black!4}
    },
    dnode sent/.style={
        dnode, fill=nodedummy!10, draw=nodedummy!50,
        rectangle split part fill={nodedummy!8, nodedummy!8, nodedummy!8}
    },
    darrow/.style={-stealth', thick, color=nodecolor!65},
    darrow prev/.style={-stealth', thick, color=nodeprev!65},
    darrow new/.style={-stealth', very thick, color=arrownew},
    darrow old/.style={-stealth', thick, color=arrowold, dashed},
    darrow new prev/.style={-stealth', very thick, color=nodeprev!80},
    dlbl/.style={font=\ttfamily\tiny\bfseries},
}


% =====================================================================
\begin{document}

\begin{frame}
	\titlepage
\end{frame}


% =============================================================
%  ROTEIRO DA AULA
% =============================================================
\begin{frame}{Roteiro da Aula}
    \begin{block}{Objetivos de Aprendizagem}
        Ao final desta aula, você será capaz de:
        \begin{enumerate}
            \item Distinguir \textbf{Tipo Abstrato de Dados} (ADT) de implementação concreta.
            \item Explicar o funcionamento e a análise amortizada de \textbf{arrays dinâmicos}.
            \item Implementar e analisar operações em \textbf{listas simplesmente ligadas}.
            \item Comparar as quatro \textbf{convenções} de organização de listas.
            \item Estender o raciocínio para \textbf{listas duplamente ligadas}.
        \end{enumerate}
    \end{block}

    \vspace{0.3cm}
    \centering
    \begin{tikzpicture}[node distance=1.8cm, >=stealth, every node/.style={font=\sffamily\scriptsize}]
        \node[draw, rounded corners, fill=nodecolor!10, minimum width=1.8cm, align=center] (a) {1. Conceitos\\Fundamentais};
        \node[draw, rounded corners, fill=nodecolor!10, minimum width=1.8cm, align=center, right of=a, xshift=1cm] (b) {2. Arrays e\\Arrays Dinâmicos};
        \node[draw, rounded corners, fill=nodecolor!10, minimum width=1.8cm, align=center, right of=b, xshift=1cm] (c) {3. Listas\\Simples};
        \node[draw, rounded corners, fill=nodecolor!10, minimum width=1.8cm, align=center, right of=c, xshift=1cm] (d) {4. Convenções\\e Variantes};
        \node[draw, rounded corners, fill=nodecolor!10, minimum width=1.8cm, align=center, right of=d, xshift=1cm] (e) {5. Listas\\Duplas};
        \draw[->] (a) -- (b);
        \draw[->] (b) -- (c);
        \draw[->] (c) -- (d);
        \draw[->] (d) -- (e);
    \end{tikzpicture}
\end{frame}


% =============================================================
%  SEÇÃO 1: CONCEITOS FUNDAMENTAIS
% =============================================================
\section{Conceitos Fundamentais}

\begin{frame}{O que é uma Estrutura de Dados?}
    \begin{block}{Definição}
        Uma estrutura de dados é a combinação de:
        \begin{enumerate}
            \item Uma \textbf{coleção de dados} organizados de forma particular.
            \item As \textbf{operações} que podem ser aplicadas a esses dados.
        \end{enumerate}
    \end{block}
    \vspace{0.3cm}
    \begin{itemize}
        \item \textbf{Analogia:} Pense em uma biblioteca. Os livros são os dados, mas as estantes e o sistema de catalogação são a \textbf{estrutura}.
        \item A escolha da estrutura determina a \textbf{complexidade} do algoritmo.
    \end{itemize}
    \vspace{0.3cm}
    \begin{block}{A Equação de Wirth}
        \centering
        Algoritmos + Estruturas de Dados = Programas
    \end{block}
\end{frame}


\begin{frame}{Tipos Abstratos de Dados (ADT)}
    \begin{columns}[T]
        \begin{column}{0.55\textwidth}
            \begin{itemize}
                \item Um \textbf{ADT} define o \textit{quê} (operações), não o \textit{como} (implementação).
                \item \textbf{Independência:} Você pode usar uma Fila sem saber se ela é feita com Array ou Lista Ligada.
                \item \textbf{Exemplos Clássicos:}
                    \begin{itemize}
                        \item Pilhas (LIFO)
                        \item Filas (FIFO)
                        \item Listas Ordenadas
                    \end{itemize}
            \end{itemize}
        \end{column}
        \begin{column}{0.43\textwidth}
            \centering
            \begin{tikzpicture}[node distance=2cm, box/.style={draw, rectangle, minimum width=2.5cm, fill=blue!5, font=\small}]
                \node[box, fill=green!10] (client) {Cliente};
                \node[box, right of=client, xshift=0.8cm] (inter) {Interface (ADT)};
                \node[box, below of=inter, yshift=0.5cm, fill=orange!10] (impl) {Implementação};

                \draw[->, thick] (client) -- node[above, font=\tiny]{\textit{usa}} (inter);
                \draw[dashed, ->, thick] (impl) -- node[right, font=\tiny]{\textit{cumpre}} (inter);
            \end{tikzpicture}
            \vspace{0.3cm}
            {\scriptsize O cliente depende apenas da interface, nunca da implementação.}
        \end{column}
    \end{columns}
\end{frame}


\begin{frame}{Estruturas Lineares vs.\ Não-Lineares}
    \centering
    \begin{columns}
        \begin{column}{0.5\textwidth}
            \begin{block}{Lineares (foco desta aula)}
                \begin{itemize}
                    \item \textbf{Arrays:} Acesso rápido via índice.
                    \item \textbf{Listas Ligadas:} Flexibilidade de inserção/remoção.
                \end{itemize}
            \end{block}
        \end{column}
        \begin{column}{0.5\textwidth}
            \begin{block}{Não-Lineares (aulas futuras)}
                \begin{itemize}
                    \item \textbf{Árvores:} Organização hierárquica.
                    \item \textbf{Grafos:} Relações arbitrárias.
                \end{itemize}
            \end{block}
        \end{column}
    \end{columns}

    \vspace{0.5cm}

    \textbf{Lista Ligada (Dinâmica)} \\
    \vspace{0.2cm}
    \begin{tikzpicture}[list/.style={rectangle split, rectangle split parts=2, draw, rectangle split horizontal}, >=stealth, node distance=1.5cm]
        \node[list] (A) {12};
        \node[list, right of=A] (B) {99};
        \node[list, right of=B] (C) {37};
        \draw[->] (A.two east) -- (B.west);
        \draw[->] (B.two east) -- (C.west);
        \draw[->] (C.two east) -- +(0.5,0) node[right] {NULL};
    \end{tikzpicture}

    \vspace{0.4cm}

    \textbf{Array (Contíguo)} \\
    \vspace{0.2cm}
    \begin{tikzpicture}
        \draw (0,0) grid (4,0.5);
        \node at (0.5, 0.25) {12};
        \node at (1.5, 0.25) {99};
        \node at (2.5, 0.25) {37};
        \node at (3.5, 0.25) {...};
        \node[scale=0.7] at (0.5, -0.3) {índice 0};
        \node[scale=0.7] at (1.5, -0.3) {índice 1};
        \node[scale=0.7] at (2.5, -0.3) {índice 2};
    \end{tikzpicture}
\end{frame}


\begin{frame}{Revisão Rápida: Referências em Java}
    \centering
    \begin{tikzpicture}[
        mem/.style={rectangle, draw, minimum width=2cm, minimum height=0.6cm},
        ref/.style={->, thick, >=stealth}
    ]
        % Primitivo
        \node[mem, fill=green!10] (int) at (0,2) {10};
        \node[left=0.2cm of int] {int n:};

        % Referência
        \node[mem, fill=blue!10] (ptr) at (0,0) {addr: 0xAF};
        \node[left=0.2cm of ptr] {Point a:};

        \node[draw, rectangle split, rectangle split parts=2, fill=yellow!10, minimum width=1.5cm] (obj) at (4,0) {
            x: 1.0 \nodepart{second} y: 1.0
        };
        \node[above=0.1cm of obj] {\tiny Endereço 0xAF};

        \draw[ref] (ptr.east) -- (obj.west);

        \draw[dashed] (-2,1) -- (6,1);
        \node at (5, 2) {\footnotesize \textbf{Direto (primitivo)}};
        \node at (5, 0.5) {\footnotesize \textbf{Indireto (referência)}};
    \end{tikzpicture}

    \vspace{0.4cm}
    \begin{columns}[T]
        \begin{column}{0.48\textwidth}
            \begin{itemize}
                \item \textbf{Primitivo:} a variável \textit{é} o valor.
                \item \textbf{Referência:} a variável \textit{aponta} para o objeto.
            \end{itemize}
        \end{column}
        \begin{column}{0.48\textwidth}
            \begin{itemize}
                \item Mais barato mover referências que copiar objetos.
                \item Base para Listas, Árvores e Grafos.
            \end{itemize}
        \end{column}
    \end{columns}
\end{frame}


% =============================================================
%  SEÇÃO 2: ARRAYS E ARRAYS DINÂMICOS
% =============================================================
\section{Arrays e Arrays Dinâmicos}

\begin{frame}{Arrays: A Base de Dados Contíguos}
    \begin{columns}[T]
        \begin{column}{0.55\textwidth}
            \begin{itemize}
                \item \textbf{Definição:} Coleção fixa de elementos do mesmo tipo.
                \item \textbf{Propriedades:}
                    \begin{itemize}
                        \item Armazenamento \textbf{contíguo} na memória.
                        \item Acesso direto via \textbf{índice}: $O(1)$.
                    \end{itemize}
                \item \textbf{Limitação:} Tamanho fixo após a criação.
            \end{itemize}
        \end{column}
        \begin{column}{0.43\textwidth}
            \centering
            \begin{tikzpicture}[node distance=0cm]
                \foreach \x in {0,1,2,3,4}
                    \node[draw, minimum width=1.2cm, minimum height=0.6cm] (m\x) at (\x*1.2,0) {\tiny Dado \x};
                \node[below=0.1cm of m0] {\tiny Índice 0};
                \node[below=0.1cm of m4] {\tiny Índice N-1};
                \draw[<-] (m0.north) -- +(0,0.3) node[above] {\tiny Endereço Base};
            \end{tikzpicture}
        \end{column}
    \end{columns}

    \vspace{0.4cm}
    \begin{block}{Analogia com a Memória}
        A memória RAM pode ser vista como um gigantesco array onde o ``índice'' é o endereço físico. O acesso \texttt{a[i]} traduz-se em apenas 3--4 instruções de CPU.
    \end{block}
\end{frame}


\begin{frame}[fragile,shrink]{Exemplo: Crivo de Eratóstenes}
    \begin{columns}[t]
        \begin{column}{0.4\textwidth}
            \small
            \textbf{Objetivo:} Encontrar primos até $N$.

            \begin{enumerate}
                \item Array \texttt{a[N]} iniciado como \texttt{true}.
                \item Para cada $i \ge 2$, se \texttt{a[i]} é primo, marca múltiplos como \texttt{false}.
            \end{enumerate}

            \begin{exampleblock}{Por que Array?}
                O acesso direto $a[k]$ em $O(1)$ permite saltar entre múltiplos eficientemente.
                Complexidade: $O(N \log \log N)$.
            \end{exampleblock}
        \end{column}

        \begin{column}{0.6\textwidth}
            \begin{lstlisting}[language=Java, basicstyle=\ttfamily\tiny, keywordstyle=\color{blue}, commentstyle=\color{gray}, frame=single]
class Primes {
    public static void main(String[] args) {
        int N = Integer.parseInt(args[0]);
        boolean[] a = new boolean[N];

        for (int i = 2; i < N; i++) a[i] = true;

        for (int i = 2; i*i  < N; i++) {
            if (a[i]) {
                for (int j = i*i; j < N; j += i)
                    a[j] = false;
            }
        }

        for (int i = 2; i < N; i++)
            if (a[i])
                System.out.print(i + " ");
        System.out.println();
    }
}
            \end{lstlisting}
        \end{column}
    \end{columns}
\end{frame}


\begin{frame}{O Problema dos Arrays Estáticos}
    \begin{alertblock}{Limitação}
        O tamanho de um array é fixo na criação. Se ele encher, não há como adicionar mais elementos.
    \end{alertblock}
    \vspace{0.3cm}
    \begin{itemize}
        \item \textbf{Solução:} O \textbf{Array Dinâmico} (como o \texttt{ArrayList} do Java) cresce automaticamente conforme necessário.
        \item \textbf{Funcionamento:}
        \begin{enumerate}
            \item Aloca um array interno com uma capacidade inicial.
            \item Quando a capacidade é atingida, cria um novo array \textbf{maior}.
            \item Copia os elementos do antigo para o novo.
        \end{enumerate}
    \end{itemize}

    \vspace{0.3cm}
    \centering
    \begin{tikzpicture}[>=stealth, scale=0.85]
        % Array antigo
        \node[font=\scriptsize\bfseries] at (-1.5, 0.25) {Antigo:};
        \foreach \x in {0,1,2,3} {
            \node[draw, minimum width=0.9cm, minimum height=0.5cm, fill=nodecolor!15] (a\x) at (\x*0.9, 0) {};
        }
        \node[font=\tiny, text=nodehigh] at (4.2, 0.25) {CHEIO!};

        % Seta
        \draw[->, very thick, color=nodeinsert] (4.5, 0) -- (5.5, 0) node[midway, above, font=\tiny] {resize};

        % Array novo
        \node[font=\scriptsize\bfseries] at (5.2, 0.25) {Novo:};
        \foreach \x in {0,...,7} {
            \pgfmathtruncatemacro{\filled}{\x < 4 ? 1 : 0}
            \ifnum\filled=1
                \node[draw, minimum width=0.9cm, minimum height=0.5cm, fill=nodecolor!15] (b\x) at (6.5+\x*0.9, 0) {};
            \else
                \node[draw, minimum width=0.9cm, minimum height=0.5cm, fill=white] (b\x) at (6.5+\x*0.9, 0) {};
            \fi
        }
        \node[font=\tiny, text=nodeinsert] at (14.0, 0.25) {espaço livre};
    \end{tikzpicture}
\end{frame}


\begin{frame}[fragile, shrink]{Implementação: Array Dinâmico}
    \begin{columns}[T]
        \begin{column}{0.5\textwidth}
            \begin{lstlisting}[language=Java, basicstyle=\ttfamily\tiny, keywordstyle=\color{blue}]
public class MeuArrayDinamico {
    private int[] data = new int[2];
    private int size = 0;

    public void add(int value) {
        if (size == data.length) {
            resize();
        }
        data[size++] = value;
    }

    private void resize() {
        int[] newData =
            new int[data.length * 2];
        for (int i = 0; i < data.length; i++)
            newData[i] = data[i];
        data = newData;
    }
}
            \end{lstlisting}
        \end{column}
        \begin{column}{0.48\textwidth}
            \begin{block}{Custo de Inserção no Fim}
                \begin{itemize}
                    \item \textbf{Caso comum:} $O(1)$ --- apenas insere.
                    \item \textbf{Caso de expansão:} $O(n)$ --- copia tudo.
                    \item \textbf{Amortizado:} $O(1)$.
                \end{itemize}
            \end{block}
            \begin{exampleblock}{Dica de Performance}
                Se você sabe o tamanho esperado, inicialize com essa capacidade:\\
		    \texttt{\scriptsize ArrayList<String> l = new ArrayList<>(10000);}
            \end{exampleblock}
        \end{column}
    \end{columns}
\end{frame}


% =============================================================
%  SEÇÃO 3: LISTAS SIMPLESMENTE LIGADAS
% =============================================================
\section{Listas Simplesmente Ligadas}

\begin{frame}[shrink]{Por que Listas Ligadas? (1/2)}
    \begin{alertblock}{Motivação}
        Arrays são ótimos para \textbf{acesso por índice}, mas péssimos para \textbf{inserção e remoção no meio} --- exigem deslocamento de todos os elementos subsequentes ($O(n)$).
    \end{alertblock}

    \vspace{0.3cm}
    \begin{itemize}
        \item \textbf{Lista ligada:} conjunto de \textbf{nós}, cada um contendo um dado e um \textbf{link} para o próximo nó.
        \item Inserção e remoção em posição conhecida: $O(1)$.
        \item \textbf{Trade-off:} perde-se o acesso aleatório.
    \end{itemize}
\end{frame}
\begin{frame}{Por que Listas Ligadas ? (2/2)}

    \centering
    \renewcommand{\arraystretch}{1.3}
    \begin{tabular}{lcc}
        \hline
        \textbf{Operação} & \textbf{Array} & \textbf{Lista Ligada} \\
        \hline
        Acesso por índice & $O(1)$ & $O(n)$ \\
        Inserção/Remoção (posição conhecida) & $O(n)$ & $O(1)$ \\
        Uso de memória & Contíguo & +1 ponteiro/nó \\
        \hline
    \end{tabular}
\end{frame}


\begin{frame}[fragile]{A Classe Node}
    \begin{columns}[T]
        \begin{column}{0.45\textwidth}
            \begin{block}{Definição}
                \begin{lstlisting}[language=Java, basicstyle=\ttfamily\small, keywordstyle=\color{blue}]
class Node {
    Object item;
    Node next;

    Node(Object v) {
        item = v;
        next = null;
    }
}
                \end{lstlisting}
            \end{block}
        \end{column}
        \begin{column}{0.53\textwidth}
            \textbf{Componentes do Nó:}
            \begin{itemize}
                \item \texttt{item}: armazena o dado.
                \item \texttt{next}: referência para o próximo \texttt{Node}.
            \end{itemize}
            \vspace{0.3cm}
            \textbf{Acesso:}
            \begin{itemize}
                \item \texttt{x.item} --- o dado do nó \texttt{x}.
                \item \texttt{x.next} --- o próximo nó.
                \item \texttt{x.next.item} --- dado do seguinte.
            \end{itemize}
            \vspace{0.2cm}
            \begin{alertblock}{\small Cuidado}
                \scriptsize Sempre verifique se \texttt{x} e \texttt{x.next} não são \texttt{null} antes do acesso, para evitar \texttt{NullPointerException}.
            \end{alertblock}
        \end{column}
    \end{columns}
\end{frame}


\begin{frame}[fragile,shrink]{Operações Fundamentais: Inserção e Remoção}
    \begin{columns}[T]
        \begin{column}{0.5\textwidth}
            \begin{block}{Remoção após \texttt{x}}
                \begin{lstlisting}[language=Java, basicstyle=\ttfamily\scriptsize]
t = x.next;
x.next = t.next;
                \end{lstlisting}
            \end{block}
        \end{column}
        \begin{column}{0.5\textwidth}
            \begin{block}{Inserção de \texttt{t} após \texttt{x}}
                \begin{lstlisting}[language=Java, basicstyle=\ttfamily\scriptsize]
t.next = x.next;
x.next = t;
                \end{lstlisting}
            \end{block}
        \end{column}
    \end{columns}
    \vspace{0.3cm}
    \centering
    \begin{tikzpicture}[scale=0.85, every node/.style={transform shape}]
        % === REMOÇÃO (esquerda) ===
        \node[step label] at (-0.3, 0.6) {Remoção:};
        \node[llnode] (ra) at (0,0) {A};
        \node[llnode high] (rx) at (1.8,0) {x};
        \node[llnode ghost] (rt) at (3.6,0) {t};
        \node[llnode] (rb) at (5.4,0) {B};
        \draw[ptr] (ra) -- (rx);
        \draw[ptr old] (rx) -- (rt);
        \draw[ptr old] (rt) -- (rb);
        \draw[ptr new] (rx) to[bend left=40] (rb);
        \node[font=\Huge, text=nodehigh] at (3.6, -0.55) {$\times$};

        % === INSERÇÃO (direita) ===
        \node[step label] at (7.5, 0.6) {Inserção:};
        \node[llnode] (ia) at (7.8,0) {A};
        \node[llnode high] (ix) at (9.6,0) {x};
        \node[llnode] (ib) at (12.6,0) {B};
        \node[llnode insert] (it) at (11.1,1.1) {t};
        \draw[ptr] (ia) -- (ix);
        \draw[ptr old] (ix) -- (ib);
        \draw[ptr new] (ix) -- (it);
        \draw[ptr new] (it) -- (ib);
    \end{tikzpicture}

    \vspace{0.2cm}
    \begin{alertblock}{\small Atenção à Ordem!}
        \scriptsize Na inserção, faça \texttt{t.next = x.next} \textbf{antes} de \texttt{x.next = t}. Inverter a ordem perde a referência ao restante da lista!
    \end{alertblock}
\end{frame}


\begin{frame}[fragile,shrink]{Inserção: Passo a Passo}
    \begin{columns}[T]
        \begin{column}{0.45\textwidth}
            \begin{block}{Código}
                \begin{lstlisting}[language=Java, basicstyle=\ttfamily\scriptsize]
t.next = x.next;
x.next = t;
                \end{lstlisting}
            \end{block}
        \end{column}
        \begin{column}{0.55\textwidth}
            \small Inserir o nó \texttt{t} \textbf{após} \texttt{x}.
        \end{column}
    \end{columns}

    \vspace{0.2cm}

    % --- Estado Inicial ---
    \begin{center}
    \textbf{Estado Inicial:}\\[0.2cm]
    \begin{tikzpicture}
        \node[llnode] (a) at (0,0) {A};
        \node[llnode high] (x) at (2,0) {x};
        \node[llnode] (b) at (4,0) {B};
        \node[llnode] (c) at (6,0) {C};
        \node[null node] (null) at (7.5,0) {null};
        \draw[ptr] (a) -- (x);
        \draw[ptr] (x) -- (b);
        \draw[ptr] (b) -- (c);
        \draw[ptr] (c) -- (null);
        \node[lbl] at (2, -0.6) {x};
        \node[lbl] at (4, -0.6) {x.next};
        \node[llnode insert] (t) at (3, 1.5) {t};
        \node[lbl green] at (3, 2.1) {novo nó};
        \draw[ptr, color=nodeinsert!60, dashed] (t) -- ++(0.8, 0) node[null node, right=-2pt] {null};
    \end{tikzpicture}
    \end{center}

    \vspace{0.1cm}\pause

    % --- Passo 1 ---
    \begin{center}
    \textbf{Passo 1:} \texttt{t.next = x.next} \hfill \textcolor{black!50}{\scriptsize (t aponta para o sucessor de x)}\\[0.2cm]
    \begin{tikzpicture}
        \node[llnode] (a) at (0,0) {A};
        \node[llnode high] (x) at (2,0) {x};
        \node[llnode] (b) at (4,0) {B};
        \node[llnode] (c) at (6,0) {C};
        \node[null node] (null) at (7.5,0) {null};
        \draw[ptr] (a) -- (x);
        \draw[ptr] (x) -- (b);
        \draw[ptr] (b) -- (c);
        \draw[ptr] (c) -- (null);
        \node[llnode insert] (t) at (3, 1.5) {t};
        \draw[ptr new] (t) -- (b);
        \node[lbl] at (2, -0.6) {x};
        \node[lbl] at (4, -0.6) {x.next = t.next};
    \end{tikzpicture}
    \end{center}

    \vspace{0.1cm}\pause

    % --- Passo 2 ---
    \begin{center}
    \textbf{Passo 2:} \texttt{x.next = t} \hfill \textcolor{black!50}{\scriptsize (x agora aponta para t)}\\[0.2cm]
    \begin{tikzpicture}
        \node[llnode] (a) at (0,0) {A};
        \node[llnode high] (x) at (2,0) {x};
        \node[llnode] (b) at (5.5,0) {B};
        \node[llnode] (c) at (7.5,0) {C};
        \node[null node] (null) at (9,0) {null};
        \draw[ptr] (a) -- (x);
        \draw[ptr] (b) -- (c);
        \draw[ptr] (c) -- (null);
        \node[llnode insert] (t) at (3.5, 0) {t};
        \draw[ptr new] (x) -- (t);
        \draw[ptr new] (t) -- (b);
        \node[lbl] at (2, -0.6) {x};
        \node[lbl green] at (3.5, -0.6) {x.next};
        \node[lbl] at (5.5, -0.6) {t.next};
    \end{tikzpicture}
    \end{center}
\end{frame}


\begin{frame}[shrink]{Comparação Visual: Lista Ligada vs.\ Array}
    \begin{columns}[T]
        % --- Coluna: Lista Ligada ---
        \begin{column}{0.48\textwidth}
            \begin{block}{\centering Lista Ligada --- $O(1)$}
            \centering
            \begin{tikzpicture}[scale=0.75, every node/.style={transform shape}]
                \node[step label] at (-0.6, 0.4) {Antes:};
                \node[llnode, minimum width=1cm] (l1) at (0,0) {1};
                \node[llnode high, minimum width=1cm] (l2) at (1.5,0) {2};
                \node[llnode, minimum width=1cm] (l3) at (3,0) {4};
                \node[llnode, minimum width=1cm] (l4) at (4.5,0) {5};
                \draw[ptr] (l1) -- (l2);
                \draw[ptr] (l2) -- (l3);
                \draw[ptr] (l3) -- (l4);
                \node[llnode insert, minimum width=1cm] (ln) at (2.25,1.3) {3};

                \node[step label] at (-0.6, -1.3) {Depois:};
                \node[llnode, minimum width=1cm] (l1b) at (0,-1.7) {1};
                \node[llnode high, minimum width=1cm] (l2b) at (1.4,-1.7) {2};
                \node[llnode insert, minimum width=1cm] (lnb) at (2.8,-1.7) {3};
                \node[llnode, minimum width=1cm] (l3b) at (4.2,-1.7) {4};
                \node[llnode, minimum width=1cm] (l4b) at (5.6,-1.7) {5};
                \draw[ptr] (l1b) -- (l2b);
                \draw[ptr new] (l2b) -- (lnb);
                \draw[ptr new] (lnb) -- (l3b);
                \draw[ptr] (l3b) -- (l4b);

                \draw[ptr new, dashed] (ln) -- (l3);
                \draw[ptr new, dashed] (l2) to[bend left=20] (ln);
            \end{tikzpicture}
            \end{block}
            \centering\small\color{nodeinsert!70!black} Apenas 2 ponteiros alterados!
        \end{column}

        % --- Coluna: Array ---
        \begin{column}{0.48\textwidth}
            \begin{block}{\centering Array --- $O(n)$}
            \centering
            \begin{tikzpicture}[scale=0.75, every node/.style={transform shape}]
                \node[step label] at (-0.6, 0.4) {Antes:};
                \foreach \val [count=\i from 0] in {1,2,4,5} {
                    \node[draw, rectangle, minimum width=1cm, minimum height=0.65cm,
                          fill=nodecolor!10, thick, font=\ttfamily\small]
                          (a\i) at (1.1*\i, 0) {\val};
                }
                \node[llnode insert, minimum width=1cm, minimum height=0.65cm]
                      (an) at (2.2, 1.3) {3};
                \draw[-Stealth, thick, color=nodeinsert!70, dashed] (an) -- (2.2, 0.45);

                \node[step label] at (-0.6, -1.3) {Depois:};
                \foreach \val [count=\i from 0] in {1,2,3,4,5} {
                    \pgfmathtruncatemacro{\filltest}{\i >= 3 ? 1 : 0}
                    \ifnum\filltest=1
                        \node[draw, rectangle, minimum width=1cm, minimum height=0.65cm,
                              fill=nodehigh!12, thick, font=\ttfamily\small]
                              (b\i) at (1.1*\i, -1.7) {\val};
                    \else
                        \ifnum\i=2
                            \node[draw, rectangle, minimum width=1cm, minimum height=0.65cm,
                                  fill=nodeinsert!15, thick, font=\ttfamily\small]
                                  (b\i) at (1.1*\i, -1.7) {\val};
                        \else
                            \node[draw, rectangle, minimum width=1cm, minimum height=0.65cm,
                                  fill=nodecolor!10, thick, font=\ttfamily\small]
                                  (b\i) at (1.1*\i, -1.7) {\val};
                        \fi
                    \fi
                }
                \draw[-Stealth, thick, color=nodehigh!70] (3.3, -0.4) -- (4.4, -1.25);
                \draw[-Stealth, thick, color=nodehigh!70] (2.2, -0.4) -- (3.3, -1.25);
                \node[step label, text=nodehigh!80] at (4.7, -0.85) {shift!};
            \end{tikzpicture}
            \end{block}
            \centering\small\color{nodehigh!70!black} Todos os elementos após devem se mover!
        \end{column}
    \end{columns}
\end{frame}


\begin{frame}{Limitações das Listas Simples}
    \begin{itemize}
        \item \textbf{Acesso por índice ($k$-ésimo item):}
        \begin{itemize}
            \item \textbf{Arrays:} Acesso direto via \texttt{a[k]} --- tempo constante.
            \item \textbf{Listas:} Requer travessia de $k$ links --- tempo linear.
        \end{itemize}

        \item \textbf{Operação ``não natural'':} Encontrar o item \textit{anterior} a um dado nó exige percorrer do início ($O(n)$).

        \item \textbf{Garbage Collection:} Em Java, ao remover um nó, o sistema automaticamente recupera a memória. Em C/C++, seria necessário liberar manualmente.
    \end{itemize}

    \vspace{0.3cm}
    \begin{exampleblock}{Lição}
        A simplicidade da inserção/remoção é a razão de ser das listas; a dificuldade de acesso aleatório é seu principal compromisso.
    \end{exampleblock}
\end{frame}


% =============================================================
%  SEÇÃO 3B: APLICAÇÃO — PROBLEMA DE JOSEPHUS
% =============================================================
\begin{frame}[shrink]{Aplicação: O Problema de Josephus}
    \begin{columns}[T]
        \begin{column}{0.52\textwidth}
            \begin{itemize}
                \item \textbf{Cenário:} $N$ pessoas formam um círculo.
                \item \textbf{Regra:} Elimina-se cada $M$-ésima pessoa, fechando o círculo após cada saída.
                \item \textbf{Objetivo:} Determinar quem será o último sobrevivente.
            \end{itemize}

            \begin{exampleblock}{Exemplo ($N=9, M=5$)}
                Eliminações: 5, 1, 7, 4, 3, 6, 9, 2.\\
                \textbf{Líder eleito: 8.}
            \end{exampleblock}

            \vspace{0.2cm}
            {\small \textbf{Por que Lista Ligada Circular?}}\\
            {\scriptsize Remoções eficientes $O(1)$ e navegação contínua pelo círculo.}
        \end{column}
        \begin{column}{0.46\textwidth}
            \centering
            \begin{tikzpicture}[scale=0.8, every node/.style={transform shape}]
                \def\N{9}
                \def\R{1.6}
                \foreach \i in {1,...,\N} {
                    \pgfmathsetmacro{\angle}{90 - (\i-1)*360/\N}
                    \node[jnode] (j\i) at (\angle:\R) {\i};
                }
                \foreach \i [evaluate=\i as \nexti using {int(mod(\i,\N)+1)}] in {1,...,\N} {
                    \draw[jarrow] (j\i) -- (j\nexti);
                }
                \node[font=\ttfamily\scriptsize\bfseries, text=nodecurr!80!black]
                    at (90-8*360/9:\R+0.45) {t};
                \draw[-stealth', thick, color=nodecurr!80]
                    (90-8*360/9:\R+0.32) -- (j9);
            \end{tikzpicture}
        \end{column}
    \end{columns}
\end{frame}


\begin{frame}[fragile,shrink]{Josephus: Implementação}
    \begin{columns}[T]
        \begin{column}{0.52\textwidth}
            \begin{lstlisting}[language=Java, basicstyle=\ttfamily\tiny,
                keywordstyle=\color{blue}]
class Josephus {
 public static void main(String[] args){
  int N = Integer.parseInt(args[0]);
  int M = Integer.parseInt(args[1]);
  Node t = new Node(1);
  t.next = t; // lista circular
  for (int i = 2; i <= N; i++) {
   t.next = new Node(i, t.next);
   t = t.next;
  }
  while (t != t.next) {
   for (int i = 1; i < M; i++)
    t = t.next; // pula M-1
   Out.print(t.next.item + " ");
   t.next = t.next.next; // remove
  }
  Out.println("\nLider: " + t.item);
 }
}
            \end{lstlisting}
        \end{column}
        \begin{column}{0.46\textwidth}
            \small
            \textbf{Construção da lista circular:}
            \begin{enumerate}
                \item Cria nó 1 apontando para si mesmo.
                \item Insere 2, 3, \ldots, $N$ logo após \texttt{t}, avançando \texttt{t}.
            \end{enumerate}

            \textbf{Eliminação:}
            \begin{enumerate}
                \item Avança \texttt{t} por $M{-}1$ posições.
                \item Remove \texttt{t.next} (o $M$-ésimo).
                \item Repete até restar um nó.
            \end{enumerate}
            \begin{block}{Complexidade}
                $O(N \cdot M)$ no total.\\
                Cada rodada percorre $M$ nós.
            \end{block}
        \end{column}
    \end{columns}
\end{frame}


\begin{frame}[shrink]{Josephus: Processo de Eliminação ($N=9$, $M=5$)}

    \centering
    \begin{tikzpicture}[scale=0.72, every node/.style={transform shape}]

        \def\R{1.35}

        % ======= RODADA 1: elimina 5 =======
        \begin{scope}[shift={(0,0)}]
            \node[font=\sffamily\scriptsize\bfseries, text=black!70] at (0, \R+0.65) {Rodada 1};
            \foreach \i in {1,...,9} {
                \pgfmathsetmacro{\angle}{90 - (\i-1)*40}
                \pgfmathsetmacro{\istarget}{\i == 5 ? 1 : 0}
                \pgfmathsetmacro{\iscurr}{\i == 4 ? 1 : 0}
                \ifnum\istarget=1
                    \node[jnode target] (r1n\i) at (\angle:\R) {\i};
                \else\ifnum\iscurr=1
                    \node[jnode curr] (r1n\i) at (\angle:\R) {\i};
                \else
                    \node[jnode] (r1n\i) at (\angle:\R) {\i};
                \fi\fi
            }
            \foreach \i [evaluate=\i as \nexti using {int(mod(\i,9)+1)}] in {1,...,9} {
                \draw[jarrow] (r1n\i) -- (r1n\nexti);
            }
            \node[font=\scriptsize\bfseries, text=nodehigh] at (0, -\R-0.5) {elimina \textbf{5}};
            \node[font=\Huge, text=nodehigh, opacity=0.6] at (r1n5) {$\times$};
        \end{scope}

        % ======= RODADA 2: elimina 1 =======
        \begin{scope}[shift={(3.8,0)}]
            \node[font=\sffamily\scriptsize\bfseries, text=black!70] at (0, \R+0.65) {Rodada 2};
            \foreach \i in {1,...,9} {
                \pgfmathsetmacro{\angle}{90 - (\i-1)*40}
                \pgfmathsetmacro{\istarget}{\i == 1 ? 1 : 0}
                \pgfmathsetmacro{\iscurr}{\i == 9 ? 1 : 0}
                \pgfmathsetmacro{\isdead}{\i == 5 ? 1 : 0}
                \ifnum\isdead=1
                    \node[jnode dead] (r2n\i) at (\angle:\R) {\i};
                \else\ifnum\istarget=1
                    \node[jnode target] (r2n\i) at (\angle:\R) {\i};
                \else\ifnum\iscurr=1
                    \node[jnode curr] (r2n\i) at (\angle:\R) {\i};
                \else
                    \node[jnode] (r2n\i) at (\angle:\R) {\i};
                \fi\fi\fi
            }
            \foreach \i/\nexti in {1/2, 2/3, 3/4, 4/6, 6/7, 7/8, 8/9, 9/1} {
                \draw[jarrow] (r2n\i) -- (r2n\nexti);
            }
            \node[font=\scriptsize\bfseries, text=nodehigh] at (0, -\R-0.5) {elimina \textbf{1}};
            \node[font=\Huge, text=nodehigh, opacity=0.6] at (r2n1) {$\times$};
        \end{scope}

        % ======= RODADA 3: elimina 7 =======
        \begin{scope}[shift={(7.6,0)}]
            \node[font=\sffamily\scriptsize\bfseries, text=black!70] at (0, \R+0.65) {Rodada 3};
            \foreach \i in {1,...,9} {
                \pgfmathsetmacro{\angle}{90 - (\i-1)*40}
                \pgfmathsetmacro{\istarget}{\i == 7 ? 1 : 0}
                \pgfmathsetmacro{\iscurr}{\i == 6 ? 1 : 0}
                \pgfmathsetmacro{\isdead}{(\i == 5 || \i == 1) ? 1 : 0}
                \ifnum\isdead=1
                    \node[jnode dead] (r3n\i) at (\angle:\R) {\i};
                \else\ifnum\istarget=1
                    \node[jnode target] (r3n\i) at (\angle:\R) {\i};
                \else\ifnum\iscurr=1
                    \node[jnode curr] (r3n\i) at (\angle:\R) {\i};
                \else
                    \node[jnode] (r3n\i) at (\angle:\R) {\i};
                \fi\fi\fi
            }
            \foreach \i/\nexti in {2/3, 3/4, 4/6, 6/7, 7/8, 8/9, 9/2} {
                \draw[jarrow] (r3n\i) -- (r3n\nexti);
            }
            \node[font=\scriptsize\bfseries, text=nodehigh] at (0, -\R-0.5) {elimina \textbf{7}};
            \node[font=\Huge, text=nodehigh, opacity=0.6] at (r3n7) {$\times$};
        \end{scope}

        % ======= RODADA 4: elimina 4 =======
        \begin{scope}[shift={(11.4,0)}]
            \node[font=\sffamily\scriptsize\bfseries, text=black!70] at (0, \R+0.65) {Rodada 4};
            \foreach \i in {1,...,9} {
                \pgfmathsetmacro{\angle}{90 - (\i-1)*40}
                \pgfmathsetmacro{\istarget}{\i == 4 ? 1 : 0}
                \pgfmathsetmacro{\iscurr}{\i == 3 ? 1 : 0}
                \pgfmathsetmacro{\isdead}{(\i==5 || \i==1 || \i==7) ? 1 : 0}
                \ifnum\isdead=1
                    \node[jnode dead] (r4n\i) at (\angle:\R) {\i};
                \else\ifnum\istarget=1
                    \node[jnode target] (r4n\i) at (\angle:\R) {\i};
                \else\ifnum\iscurr=1
                    \node[jnode curr] (r4n\i) at (\angle:\R) {\i};
                \else
                    \node[jnode] (r4n\i) at (\angle:\R) {\i};
                \fi\fi\fi
            }
            \foreach \i/\nexti in {2/3, 3/4, 4/6, 6/8, 8/9, 9/2} {
                \draw[jarrow] (r4n\i) -- (r4n\nexti);
            }
            \node[font=\scriptsize\bfseries, text=nodehigh] at (0, -\R-0.5) {elimina \textbf{4}};
            \node[font=\Huge, text=nodehigh, opacity=0.6] at (r4n4) {$\times$};
        \end{scope}

    \end{tikzpicture}

    \vspace{0.3cm}

    \begin{tikzpicture}[scale=0.72, every node/.style={transform shape}]

        \def\R{1.35}

        % ======= RODADA 5: elimina 3 =======
        \begin{scope}[shift={(0,0)}]
            \node[font=\sffamily\scriptsize\bfseries, text=black!70] at (0, \R+0.65) {Rodada 5};
            \foreach \i in {1,...,9} {
                \pgfmathsetmacro{\angle}{90 - (\i-1)*40}
                \pgfmathsetmacro{\istarget}{\i == 3 ? 1 : 0}
                \pgfmathsetmacro{\iscurr}{\i == 2 ? 1 : 0}
                \pgfmathsetmacro{\isdead}{(\i==5||\i==1||\i==7||\i==4) ? 1 : 0}
                \ifnum\isdead=1
                    \node[jnode dead] (r5n\i) at (\angle:\R) {\i};
                \else\ifnum\istarget=1
                    \node[jnode target] (r5n\i) at (\angle:\R) {\i};
                \else\ifnum\iscurr=1
                    \node[jnode curr] (r5n\i) at (\angle:\R) {\i};
                \else
                    \node[jnode] (r5n\i) at (\angle:\R) {\i};
                \fi\fi\fi
            }
            \foreach \i/\nexti in {2/3, 3/6, 6/8, 8/9, 9/2} {
                \draw[jarrow] (r5n\i) -- (r5n\nexti);
            }
            \node[font=\scriptsize\bfseries, text=nodehigh] at (0, -\R-0.5) {elimina \textbf{3}};
            \node[font=\Huge, text=nodehigh, opacity=0.6] at (r5n3) {$\times$};
        \end{scope}

        % ======= RODADA 6: elimina 6 =======
        \begin{scope}[shift={(3.8,0)}]
            \node[font=\sffamily\scriptsize\bfseries, text=black!70] at (0, \R+0.65) {Rodada 6};
            \foreach \i in {1,...,9} {
                \pgfmathsetmacro{\angle}{90 - (\i-1)*40}
                \pgfmathsetmacro{\istarget}{\i == 6 ? 1 : 0}
                \pgfmathsetmacro{\iscurr}{\i == 9 ? 1 : 0}
                \pgfmathsetmacro{\isdead}{(\i==5||\i==1||\i==7||\i==4||\i==3) ? 1 : 0}
                \ifnum\isdead=1
                    \node[jnode dead] (r6n\i) at (\angle:\R) {\i};
                \else\ifnum\istarget=1
                    \node[jnode target] (r6n\i) at (\angle:\R) {\i};
                \else\ifnum\iscurr=1
                    \node[jnode curr] (r6n\i) at (\angle:\R) {\i};
                \else
                    \node[jnode] (r6n\i) at (\angle:\R) {\i};
                \fi\fi\fi
            }
            \foreach \i/\nexti in {2/6, 6/8, 8/9, 9/2} {
                \draw[jarrow] (r6n\i) -- (r6n\nexti);
            }
            \node[font=\scriptsize\bfseries, text=nodehigh] at (0, -\R-0.5) {elimina \textbf{6}};
            \node[font=\Huge, text=nodehigh, opacity=0.6] at (r6n6) {$\times$};
        \end{scope}

        % ======= RODADA 7: elimina 9 =======
        \begin{scope}[shift={(7.6,0)}]
            \node[font=\sffamily\scriptsize\bfseries, text=black!70] at (0, \R+0.65) {Rodada 7};
            \foreach \i in {1,...,9} {
                \pgfmathsetmacro{\angle}{90 - (\i-1)*40}
                \pgfmathsetmacro{\istarget}{\i == 9 ? 1 : 0}
                \pgfmathsetmacro{\iscurr}{\i == 8 ? 1 : 0}
                \pgfmathsetmacro{\isdead}{(\i==5||\i==1||\i==7||\i==4||\i==3||\i==6) ? 1 : 0}
                \ifnum\isdead=1
                    \node[jnode dead] (r7n\i) at (\angle:\R) {\i};
                \else\ifnum\istarget=1
                    \node[jnode target] (r7n\i) at (\angle:\R) {\i};
                \else\ifnum\iscurr=1
                    \node[jnode curr] (r7n\i) at (\angle:\R) {\i};
                \else
                    \node[jnode] (r7n\i) at (\angle:\R) {\i};
                \fi\fi\fi
            }
            \foreach \i/\nexti in {2/8, 8/9, 9/2} {
                \draw[jarrow] (r7n\i) -- (r7n\nexti);
            }
            \node[font=\scriptsize\bfseries, text=nodehigh] at (0, -\R-0.5) {elimina \textbf{9}};
            \node[font=\Huge, text=nodehigh, opacity=0.6] at (r7n9) {$\times$};
        \end{scope}

        % ======= RODADA 8: elimina 2, sobra 8 =======
        \begin{scope}[shift={(11.4,0)}]
            \node[font=\sffamily\scriptsize\bfseries, text=black!70] at (0, \R+0.65) {Rodada 8};
            \foreach \i in {1,...,9} {
                \pgfmathsetmacro{\angle}{90 - (\i-1)*40}
                \pgfmathsetmacro{\istarget}{\i == 2 ? 1 : 0}
                \pgfmathsetmacro{\iscurr}{\i == 8 ? 1 : 0}
                \pgfmathsetmacro{\isdead}{(\i==5||\i==1||\i==7||\i==4||\i==3||\i==6||\i==9) ? 1 : 0}
                \ifnum\isdead=1
                    \node[jnode dead] (r8n\i) at (\angle:\R) {\i};
                \else\ifnum\istarget=1
                    \node[jnode target] (r8n\i) at (\angle:\R) {\i};
                \else\ifnum\iscurr=1
                    \node[jnode curr] (r8n\i) at (\angle:\R) {\i};
                \else
                    \node[jnode] (r8n\i) at (\angle:\R) {\i};
                \fi\fi\fi
            }
            \foreach \i/\nexti in {2/8, 8/2} {
                \draw[jarrow] (r8n\i) -- (r8n\nexti);
            }
            \node[font=\scriptsize\bfseries, text=nodehigh] at (0, -\R-0.5) {elimina \textbf{2}};
            \node[font=\Huge, text=nodehigh, opacity=0.6] at (r8n2) {$\times$};
            \draw[very thick, color=nodeinsert!80] (r8n8) circle (0.6cm);
        \end{scope}

    \end{tikzpicture}

    \vspace{0.2cm}
    \centering
    {\small Ordem: \textcolor{nodehigh}{5, 1, 7, 4, 3, 6, 9, 2} \qquad
    Líder: \textcolor{nodeinsert}{\textbf{8}}}
\end{frame}


\begin{frame}{Josephus: Lição de Design}
    \begin{table}[]
        \small
        \begin{tabular}{|l|l|l|}
        \hline
        \textbf{Algoritmo} & \textbf{Estrutura Ideal} & \textbf{Motivo} \\ \hline
        Crivo de Eratóstenes & \textbf{Array} & Acesso aleatório rápido via índice. \\ \hline
        Problema de Josephus & \textbf{Lista Circular} & Remoção eficiente de itens. \\ \hline
        \end{tabular}
    \end{table}

    \vspace{0.3cm}
    \begin{alertblock}{Lição Fundamental}
        O uso de uma lista para o Crivo seria caro (falta de acesso direto). O uso de um array para Josephus seria caro (deslocamento na remoção). O design eficiente nasce da \textbf{interação} entre estrutura e algoritmo.
    \end{alertblock}
\end{frame}


% =============================================================
%  SEÇÃO 4: CONVENÇÕES DE ORGANIZAÇÃO
% =============================================================
\section{Convenções de Organização de Listas}

\begin{frame}[shrink]{As 4 Convenções de Organização}
    \centering
    \begin{tikzpicture}[scale=0.82, every node/.style={transform shape}]

        % ====== 1. Circular ======
        \begin{scope}[shift={(0, 0)}]
            \node[conv title, anchor=west] at (-2.2, 0.9) {1. Circular, nunca vazia};
            \node[lnode] (c1) at (0, 0) {A \nodepart{two} \phantom{.}};
            \node[lnode] (c2) at (2, 0) {B \nodepart{two} \phantom{.}};
            \node[lnode] (c3) at (4, 0) {C \nodepart{two} \phantom{.}};
            \draw[larrow] (c1.two east) -- (c2);
            \draw[larrow] (c2.two east) -- (c3);
            \draw[larrow] (c3.two east) to[out=0, in=0, looseness=1.6] ($(c3.east)+(0.3,0.7)$) -- ($(c1.west)+(-0.3,0.7)$) to[out=180,in=180,looseness=0.8] (c1.west);
            \node[head ptr] at (-0.05, -0.6) {head};
            \draw[harrow] (-0.05, -0.45) -- (c1.south);
        \end{scope}

        % ====== 2. Head/null ======
        \begin{scope}[shift={(7.5, 0)}]
            \node[conv title, anchor=west] at (-2.2, 0.9) {2. Ref.\ head, cauda null};
            \node[lnode] (h1) at (0, 0) {A \nodepart{two} \phantom{.}};
            \node[lnode] (h2) at (2, 0) {B \nodepart{two} \phantom{.}};
            \node[lnode] (h3) at (4, 0) {C \nodepart{two} \phantom{.}};
            \node[null text] (hn) at (5.7, 0) {null};
            \draw[larrow] (h1.two east) -- (h2);
            \draw[larrow] (h2.two east) -- (h3);
            \draw[larrow] (h3.two east) -- (hn);
            \node[head ptr] at (-0.05, -0.6) {head};
            \draw[harrow] (-0.05, -0.45) -- (h1.south);
        \end{scope}

        % ====== 3. Dummy head ======
        \begin{scope}[shift={(0, -2.5)}]
            \node[conv title, anchor=west] at (-2.2, 0.9) {3. Nó sentinela head, cauda null};
            \node[lnode sent] (d1) at (0, 0) {\textcolor{nodedummy!70}{$\star$} \nodepart{two} \phantom{.}};
            \node[lnode] (d2) at (2, 0) {A \nodepart{two} \phantom{.}};
            \node[lnode] (d3) at (4, 0) {B \nodepart{two} \phantom{.}};
            \node[null text] (dn) at (5.7, 0) {null};
            \draw[larrow] (d1.two east) -- (d2);
            \draw[larrow] (d2.two east) -- (d3);
            \draw[larrow] (d3.two east) -- (dn);
            \node[head ptr] at (-0.05, -0.6) {head};
            \draw[harrow] (-0.05, -0.45) -- (d1.south);
            \node[conv subtitle] at (0, -0.9) {(dummy)};
        \end{scope}

        % ====== 4. Dummy head + tail ======
        \begin{scope}[shift={(7.5, -2.5)}]
            \node[conv title, anchor=west] at (-2.2, 0.9) {4. Sentinelas head e tail};
            \node[lnode sent] (s1) at (0, 0) {\textcolor{nodedummy!70}{$\star$} \nodepart{two} \phantom{.}};
            \node[lnode] (s2) at (2, 0) {A \nodepart{two} \phantom{.}};
            \node[lnode] (s3) at (4, 0) {B \nodepart{two} \phantom{.}};
            \node[lnode sent] (sz) at (6, 0) {\textcolor{nodedummy!70}{$\star$} \nodepart{two} \phantom{.}};
            \draw[larrow] (s1.two east) -- (s2);
            \draw[larrow] (s2.two east) -- (s3);
            \draw[larrow] (s3.two east) -- (sz);
            \draw[larrow] (sz.two east) to[out=0, in=0, looseness=4] (sz.east);
            \node[head ptr] at (-0.05, -0.6) {head};
            \draw[harrow] (-0.05, -0.45) -- (s1.south);
            \node[head ptr, text=nodedummy!80!black] at (6, -0.6) {z};
            \draw[-stealth', very thick, color=nodedummy!70] (6, -0.45) -- (sz.south);
            \node[conv subtitle] at (0, -0.9) {(dummy)};
            \node[conv subtitle] at (6, -0.9) {(dummy)};
        \end{scope}

    \end{tikzpicture}
\end{frame}


\begin{frame}{Comparação das 4 Convenções}

    \centering
    \scriptsize
    \renewcommand{\arraystretch}{1.4}
    \begin{tabular}{@{} l c c c c @{}}
        \hline
        \textbf{Propriedade}
            & \textbf{1. Circular}
            & \textbf{2. Head/Null}
            & \textbf{3. Dummy Head}
            & \textbf{4. Dummy H+T} \\
        \hline
        Lista vazia?
            & Não
            & Sim
            & Sim
            & Sim \\
        Caso especial na inserção
            & Não
            & \textcolor{nodehigh}{Sim}
            & Não
            & Não \\
        Testa \texttt{null} na travessia
            & Não
            & Sim
            & Sim
            & Não \\
        Memória extra
            & Nenhuma
            & Nenhuma
            & +1 nó
            & +2 nós \\
        Risco de \texttt{NullPtrExc}
            & Não
            & \textcolor{nodehigh}{Sim}
            & \textcolor{nodehigh}{Sim}
            & Não \\
        Uso típico
            & \parbox{2cm}{\centering Josephus,\\round-robin}
            & \parbox{2cm}{\centering Uso geral,\\pilhas/filas}
            & \parbox{2cm}{\centering Código\\uniforme}
            & \parbox{2cm}{\centering Busca\\sequencial} \\
        \hline
    \end{tabular}

    \vspace{0.4cm}

    \begin{tikzpicture}[scale=0.6, every node/.style={transform shape}]

        \begin{scope}[shift={(0,0)}]
            \node[font=\sffamily\tiny\bfseries, text=black!60] at (1, 1) {1. Circular};
            \node[lnode, minimum width=0.8cm] (c1) at (0,0) {\tiny A};
            \node[lnode, minimum width=0.8cm] (c2) at (1.3,0) {\tiny B};
            \node[lnode, minimum width=0.8cm] (c3) at (2.6,0) {\tiny C};
            \draw[larrow, thin] (c1) -- (c2);
            \draw[larrow, thin] (c2) -- (c3);
            \draw[larrow, thin, color=nodehigh!50] (c3.east) to[out=0,in=0,looseness=1.5]
                ($(c3.east)+(0.15,0.5)$) -- ($(c1.west)+(-0.15,0.5)$)
                to[out=180,in=180,looseness=0.8] (c1.west);
        \end{scope}

        \begin{scope}[shift={(4.5,0)}]
            \node[font=\sffamily\tiny\bfseries, text=black!60] at (1, 1) {2. Head/null};
            \node[lnode, minimum width=0.8cm] (h1) at (0,0) {\tiny A};
            \node[lnode, minimum width=0.8cm] (h2) at (1.3,0) {\tiny B};
            \node[null text] (hn) at (2.5,0) {\tiny null};
            \draw[larrow, thin] (h1) -- (h2);
            \draw[larrow, thin] (h2) -- (hn);
        \end{scope}

        \begin{scope}[shift={(8.5,0)}]
            \node[font=\sffamily\tiny\bfseries, text=black!60] at (1.2, 1) {3. Dummy head};
            \node[lnode sent, minimum width=0.8cm] (d0) at (0,0) {\tiny $\star$};
            \node[lnode, minimum width=0.8cm] (d1) at (1.3,0) {\tiny A};
            \node[lnode, minimum width=0.8cm] (d2) at (2.6,0) {\tiny B};
            \node[null text] (dn) at (3.8,0) {\tiny null};
            \draw[larrow, thin] (d0) -- (d1);
            \draw[larrow, thin] (d1) -- (d2);
            \draw[larrow, thin] (d2) -- (dn);
        \end{scope}

        \begin{scope}[shift={(13.5,0)}]
            \node[font=\sffamily\tiny\bfseries, text=black!60] at (1.2, 1) {4. Dummy H+T};
            \node[lnode sent, minimum width=0.8cm] (s0) at (0,0) {\tiny $\star$};
            \node[lnode, minimum width=0.8cm] (s1) at (1.3,0) {\tiny A};
            \node[lnode sent, minimum width=0.8cm] (s2) at (2.6,0) {\tiny $\star$};
            \draw[larrow, thin] (s0) -- (s1);
            \draw[larrow, thin] (s1) -- (s2);
            \draw[larrow, thin, color=nodedummy!40] (s2.east) to[out=0,in=0,looseness=5] (s2.east);
        \end{scope}

    \end{tikzpicture}
\end{frame}


% =============================================================
%  SEÇÃO 5: LISTAS DUPLAMENTE LIGADAS
% =============================================================
\section{Listas Duplamente Ligadas}

\begin{frame}{Listas Duplamente Ligadas: Estrutura}

    \begin{columns}[T]
        \begin{column}{0.55\textwidth}
            \small
            Cada nó mantém \textbf{dois} ponteiros:
            \begin{itemize}
                \item \texttt{prev} --- nó anterior
                \item \texttt{next} --- nó seguinte
            \end{itemize}

            \vspace{0.2cm}
            Permite percorrer a lista em \textbf{ambas as direções} e remover um nó conhecendo \textbf{apenas} o seu ponteiro.

            \vspace{0.3cm}
            \textbf{Invariante fundamental:}\\[0.1cm]
            {\centering
            \colorbox{nodecolor!8}{\texttt{x == x.next.prev == x.prev.next}}\par}
        \end{column}
        \begin{column}{0.43\textwidth}
            \centering
            \begin{tikzpicture}[scale=0.9, every node/.style={transform shape}]
                \node[draw, rectangle split, rectangle split parts=3,
                      rectangle split horizontal, minimum height=1cm,
                      minimum width=3.5cm, thick, rounded corners=2pt,
                      rectangle split part fill={nodeprev!15, nodecolor!12, nodecolor!15},
                      font=\ttfamily\small]
                      (box) at (0,0)
                      {\textcolor{nodeprev!70}{prev}
                       \nodepart{two} \textbf{item}
                       \nodepart{three} \textcolor{nodecolor!70}{next}};

                \draw[-stealth', very thick, color=nodeprev!60]
                      (box.one west)+(-0.1,0) -- ++(-1.0,0)
                      node[left, font=\scriptsize\itshape, text=nodeprev!70] {anterior};
                \draw[-stealth', very thick, color=nodecolor!60]
                      (box.three east)+(0.1,0) -- ++(1.0,0)
                      node[right, font=\scriptsize\itshape, text=nodecolor!70] {seguinte};
            \end{tikzpicture}
        \end{column}
    \end{columns}

    \vspace{0.3cm}
    \centering
    \begin{tikzpicture}[scale=0.82, every node/.style={transform shape}]
        \node[dnode] (a) at (0,0)   {\phantom{.} \nodepart{two} A \nodepart{three} \phantom{.}};
        \node[dnode] (b) at (2.8,0) {\phantom{.} \nodepart{two} B \nodepart{three} \phantom{.}};
        \node[dnode] (c) at (5.6,0) {\phantom{.} \nodepart{two} C \nodepart{three} \phantom{.}};
        \node[dnode] (d) at (8.4,0) {\phantom{.} \nodepart{two} D \nodepart{three} \phantom{.}};

        \draw[darrow] (a.three east) to[out=15, in=165] (b.one west);
        \draw[darrow] (b.three east) to[out=15, in=165] (c.one west);
        \draw[darrow] (c.three east) to[out=15, in=165] (d.one west);

        \draw[darrow prev] (b.one west) to[out=195, in=345] (a.three east);
        \draw[darrow prev] (c.one west) to[out=195, in=345] (b.three east);
        \draw[darrow prev] (d.one west) to[out=195, in=345] (c.three east);

        \node[font=\sffamily\tiny\itshape, text=nodecolor!70] at (1.4, 0.55) {next};
        \node[font=\sffamily\tiny\itshape, text=nodeprev!70] at (1.4, -0.55) {prev};

        \node[font=\ttfamily\tiny\bfseries, text=black!30] (nl) at (-1.5, 0) {null};
        \node[font=\ttfamily\tiny\bfseries, text=black!30] (nr) at (9.9, 0) {null};
        \draw[darrow prev, color=black!25] (a.one west) -- (nl);
        \draw[darrow, color=black!25] (d.three east) -- (nr);
    \end{tikzpicture}

    \vspace{0.3cm}
    \begin{alertblock}{\small Custo}
        \scriptsize Dobro de espaço para ponteiros e dobro de manipulações por operação.
        Usar apenas quando travessia bidirecional ou remoção direta for necessária.
    \end{alertblock}
\end{frame}


\begin{frame}[fragile]{Lista Duplamente Ligada: Remoção}

    \begin{columns}[T]
        \begin{column}{0.42\textwidth}
            \begin{block}{Código}
                \begin{lstlisting}[language=Java, basicstyle=\ttfamily\scriptsize,
                    keywordstyle=\color{blue}]
// Remover o no x
x.prev.next = x.next;
x.next.prev = x.prev;
                \end{lstlisting}
            \end{block}
            \vspace{0.1cm}
            \small
            \textbf{Diferença crucial:}\\
            {\scriptsize Na lista simples, precisamos do nó \emph{anterior} a \texttt{x}. Na dupla, \texttt{x} já contém \texttt{x.prev}.}
        \end{column}
        \begin{column}{0.56\textwidth}
            \centering\small
            Basta conhecer \texttt{x} para removê-lo!\\
            Nenhuma busca pelo predecessor.
        \end{column}
    \end{columns}

    \vspace{0.25cm}

    \centering
    \textbf{Antes:}\\[0.15cm]
    \begin{tikzpicture}[scale=0.82, every node/.style={transform shape}]
        \node[dnode] (a) at (0,0)   {\phantom{.} \nodepart{two} A \nodepart{three} \phantom{.}};
        \node[dnode hi] (x) at (3,0) {\phantom{.} \nodepart{two} x \nodepart{three} \phantom{.}};
        \node[dnode] (c) at (6,0)   {\phantom{.} \nodepart{two} C \nodepart{three} \phantom{.}};

        \draw[darrow] (a.three east) to[out=15, in=165] (x.one west);
        \draw[darrow] (x.three east) to[out=15, in=165] (c.one west);
        \draw[darrow prev] (x.one west) to[out=195, in=345] (a.three east);
        \draw[darrow prev] (c.one west) to[out=195, in=345] (x.three east);

        \node[dlbl, text=nodehigh!80!black] at (3, -0.7) {x};
        \node[dlbl, text=black!50] at (0, -0.7) {x.prev};
        \node[dlbl, text=black!50] at (6, -0.7) {x.next};
    \end{tikzpicture}

    \vspace{0.15cm}\pause

    \textbf{Depois:} \texttt{x.prev.next = x.next; x.next.prev = x.prev}\\[0.15cm]
    \begin{tikzpicture}[scale=0.82, every node/.style={transform shape}]
        \node[dnode] (a) at (0,0)   {\phantom{.} \nodepart{two} A \nodepart{three} \phantom{.}};
        \node[dnode ghost] (x) at (3,0) {\phantom{.} \nodepart{two} x \nodepart{three} \phantom{.}};
        \node[dnode] (c) at (6,0)   {\phantom{.} \nodepart{two} C \nodepart{three} \phantom{.}};

        \draw[darrow new] (a.three east) to[out=25, in=155] (c.one west);
        \draw[darrow new prev] (c.one west) to[out=205, in=335] (a.three east);

        \node[font=\Large, text=nodehigh, opacity=0.6] at (3, 0) {$\times$};
        \node[font=\sffamily\tiny\itshape, text=black!40] at (3, -0.7) {desconectado};
    \end{tikzpicture}
\end{frame}


\begin{frame}[fragile]{Lista Duplamente Ligada: Inserção}

    \begin{columns}[T]
        \begin{column}{0.48\textwidth}
            \begin{block}{\small Inserir \texttt{t} \textbf{após} \texttt{x}}
                \begin{lstlisting}[language=Java, basicstyle=\ttfamily\tiny,
                    keywordstyle=\color{blue}]
t.next = x.next;
t.prev = x;
x.next.prev = t;
x.next = t;
                \end{lstlisting}
            \end{block}
        \end{column}
        \begin{column}{0.48\textwidth}
            \begin{block}{\small Inserir \texttt{t} \textbf{antes} de \texttt{x}}
                \begin{lstlisting}[language=Java, basicstyle=\ttfamily\tiny,
                    keywordstyle=\color{blue}]
t.prev = x.prev;
t.next = x;
x.prev.next = t;
x.prev = t;
                \end{lstlisting}
            \end{block}
        \end{column}
    \end{columns}

    \vspace{0.2cm}

    \centering
    {\small\bfseries Inserir \texttt{t} após \texttt{x}:}

    \vspace{0.1cm}
    \begin{tikzpicture}[scale=0.75, every node/.style={transform shape}]
        \node[font=\sffamily\scriptsize\itshape, text=black!50] at (-1.8, 0.6) {Antes:};
        \node[dnode] (a1) at (0,0.6)   {\phantom{.} \nodepart{two} A \nodepart{three} \phantom{.}};
        \node[dnode hi] (x1) at (3,0.6) {\phantom{.} \nodepart{two} x \nodepart{three} \phantom{.}};
        \node[dnode] (b1) at (6,0.6)   {\phantom{.} \nodepart{two} B \nodepart{three} \phantom{.}};
        \draw[darrow] (a1.three east) to[out=12,in=168] (x1.one west);
        \draw[darrow] (x1.three east) to[out=12,in=168] (b1.one west);
        \draw[darrow prev] (x1.one west) to[out=198,in=342] (a1.three east);
        \draw[darrow prev] (b1.one west) to[out=198,in=342] (x1.three east);

        \node[dnode new] (t1) at (4.5, 2.0) {\phantom{.} \nodepart{two} t \nodepart{three} \phantom{.}};
        \node[dlbl, text=nodeinsert!80!black] at (4.5, 2.65) {novo};

        \node[font=\sffamily\scriptsize\itshape, text=black!50] at (-1.8, -1.4) {Depois:};
        \node[dnode] (a2) at (0,-1.4)   {\phantom{.} \nodepart{two} A \nodepart{three} \phantom{.}};
        \node[dnode hi] (x2) at (2.7,-1.4) {\phantom{.} \nodepart{two} x \nodepart{three} \phantom{.}};
        \node[dnode new] (t2) at (5.4,-1.4) {\phantom{.} \nodepart{two} t \nodepart{three} \phantom{.}};
        \node[dnode] (b2) at (8.1,-1.4) {\phantom{.} \nodepart{two} B \nodepart{three} \phantom{.}};

        \draw[darrow] (a2.three east) to[out=12,in=168] (x2.one west);
        \draw[darrow prev] (x2.one west) to[out=198,in=342] (a2.three east);
        \draw[darrow new] (x2.three east) to[out=12,in=168] (t2.one west);
        \draw[darrow new prev] (t2.one west) to[out=198,in=342] (x2.three east);
        \draw[darrow new] (t2.three east) to[out=12,in=168] (b2.one west);
        \draw[darrow new prev] (b2.one west) to[out=198,in=342] (t2.three east);
    \end{tikzpicture}
\end{frame}


\begin{frame}{Comparação: Lista Simples vs.\ Duplamente Ligada}
   \renewcommand{\arraystretch}{1.3}
    \centering
    \begin{tabular}{@{} l c c @{}}
        \hline
        \textbf{Operação} & \textbf{Lista Simples} & \textbf{Lista Dupla} \\
        \hline
        Inserir após \texttt{x} & $O(1)$ --- 2 ponteiros & $O(1)$ --- 4 ponteiros \\
        Inserir antes de \texttt{x} & $O(n)$ --- buscar predecessor & $O(1)$ --- 4 ponteiros \\
        Remover \texttt{x} (dado \texttt{x}) & $O(n)$ --- buscar predecessor & $O(1)$ --- 2 ponteiros \\
        Remover após \texttt{x} & $O(1)$ --- 1 ponteiro & $O(1)$ --- 2 ponteiros \\
        Espaço por nó & 1 ponteiro & 2 ponteiros \\
        \hline
    \end{tabular}

    \vspace{0.4cm}

    \begin{columns}[T]
        \begin{column}{0.48\textwidth}
            \centering
            \begin{tikzpicture}
                \node[draw, rounded corners=4pt, fill=nodehigh!6, draw=nodehigh!40,
                      thick, font=\sffamily\tiny, text width=4.5cm, align=left,
                      inner sep=4pt]
                {
                    \textbf{Quando usar lista simples:}\\
                    Inserção/remoção sempre no início ou após um nó conhecido.
                    Economia de memória.
                };
            \end{tikzpicture}
        \end{column}
        \begin{column}{0.48\textwidth}
            \centering
            \begin{tikzpicture}
                \node[draw, rounded corners=4pt, fill=nodeinsert!6, draw=nodeinsert!40,
                      thick, font=\sffamily\tiny, text width=4.5cm, align=left,
                      inner sep=4pt]
                {
                    \textbf{Quando usar lista dupla:}\\
                    Remoção direta de nó, travessia bidirecional,
                    nós compartilhados entre estruturas (ex: LRU cache).
                };
            \end{tikzpicture}
        \end{column}
    \end{columns}
\end{frame}


% =============================================================
%  RECAPITULAÇÃO
% =============================================================
\section{Recapitulação}

\begin{frame}{Recapitulação: O que Aprendemos Hoje}
    \centering
    \renewcommand{\arraystretch}{1.4}
    \begin{tabular}{@{} l c c c @{}}
        \hline
        \textbf{Operação} & \textbf{Array Estático} & \textbf{Array Dinâmico} & \textbf{Lista Ligada} \\
        \hline
        Acesso por índice & $O(1)$ & $O(1)$ & $O(n)$ \\
        Inserção no fim & --- & $O(1)$ amortizado & $O(1)$ \\
        Inserção no meio & $O(n)$ & $O(n)$ & $O(1)$\textsuperscript{*} \\
        Remoção (posição conhecida) & $O(n)$ & $O(n)$ & $O(1)$\textsuperscript{*} \\
        Memória & Contígua, fixa & Contígua, cresce & Fragmentada \\
        \hline
    \end{tabular}

    {\scriptsize \textsuperscript{*}Se já temos a referência ao nó.}

    \vspace{0.4cm}
    \begin{block}{Mensagem Central}
        Não existe estrutura ``melhor'' em absoluto. A escolha depende do \textbf{padrão de acesso} do algoritmo que a utiliza:
        \begin{itemize}
            \item Acesso aleatório frequente $\Rightarrow$ \textbf{Array}.
            \item Inserções/remoções frequentes em posições variadas $\Rightarrow$ \textbf{Lista Ligada}.
        \end{itemize}
    \end{block}
\end{frame}


\begin{frame}{Próximos Passos}
    \begin{itemize}
        \item \textbf{Próxima aula:} Pilhas e Filas --- ADTs construídos sobre arrays e listas.
        \item \textbf{Exercícios sugeridos:}
            \begin{enumerate}
                \item Implemente um \texttt{ArrayList} simplificado com operação \texttt{remove(int index)}.
                \item Resolva o problema de Josephus para diferentes valores de $N$ e $M$.
                \item Implemente uma lista duplamente ligada com sentinelas e operações de inserção/remoção.
            \end{enumerate}
        \item \textbf{Leitura:} Sedgewick, capítulos sobre estruturas elementares.
    \end{itemize}
\end{frame}


\end{document}
